\documentclass[11pt, a4paper]{article}

\usepackage{amsmath, amssymb, amsfonts}
\usepackage{graphicx}
\usepackage{geometry}
\usepackage{hyperref}
\usepackage{tikz}
\usepackage{booktabs}
\usepackage{xcolor}
\usepackage{bm}
\usepackage{float}
\usetikzlibrary{arrows.meta, positioning, fit, calc, shapes.geometric}

\geometry{margin=2.5cm}

\title{%
  \textbf{LQR with Optimal Feedback Control (OFC)} \\[0.5em]
  \large Mode: \texttt{lqr-applied-to-cart-manip-following-cart} \\[0.2em]
  \normalsize 2D Cart-Pendulum System with Muscle Dynamics and Manipulator Tracking
}
\author{Research Note}
\date{\today}

\begin{document}
\maketitle
\tableofcontents
\newpage

% ============================================================================
\section{Overview}
% ============================================================================

This note describes the formulation and implementation of the
\texttt{lqr-applied-to-cart-manip-following-cart} simulation mode.
The system consists of two coupled sub-systems:

\begin{enumerate}
  \item \textbf{Cart-Pendulum (4~DOF):} A 2D cart moving in the $XY$-plane
        carrying a 3D gimbal pendulum. It is controlled by a biologically
        inspired Optimal Feedback Control (OFC) loop using Muscle Dynamics,
        a Zero-Force-Trajectory (ZFT) Reference Mass, Impedance Force, and a
        Finite-Horizon LQR.
  \item \textbf{Cup Manipulator (2~DOF):} A 2-link revolute manipulator whose
        end-effector (EE) is driven to track the cart's 2D position via online
        Inverse Kinematics (IK) and a Joint-Space Computed-Torque Controller.
\end{enumerate}

The two sub-systems share the same Drake \texttt{MultibodyPlant} but have
\emph{separate model instances} so their states and actuations can be queried
and applied independently.

% ============================================================================
\section{State and Coordinate Definitions}
% ============================================================================

\subsection{Cart-Pendulum State (8D)}

\[
\bm{s}_{\text{cart}} =
\underbrace{[x,\; y]^{\top}}_{\text{cart position}}
\oplus
\underbrace{[\alpha,\; \beta]^{\top}}_{\text{pend. angles}}
\oplus
\underbrace{[\dot{x},\; \dot{y}]^{\top}}_{\text{cart velocity}}
\oplus
\underbrace{[\dot{\alpha},\; \dot{\beta}]^{\top}}_{\text{pend. ang. vel.}}
\in \mathbb{R}^{8}
\]

where $\alpha$ is the pendulum pitch and $\beta$ is the roll (gimbal angles).

\subsection{Full OFC State (14D)}

The LQR operates on the extended state:

\[
\bm{x} =
\begin{bmatrix}
x,\; y,\; \alpha,\; \beta,\;
\dot{x},\; \dot{y},\; \dot{\alpha},\; \dot{\beta},\;
F_x,\; F_y,\;
x_{\mathrm{ref}},\; y_{\mathrm{ref}},\;
\dot{x}_{\mathrm{ref}},\; \dot{y}_{\mathrm{ref}}
\end{bmatrix}^{\top} \in \mathbb{R}^{14}
\]

\subsection{Manipulator State (4D)}

\[
\bm{s}_{\text{manip}} =
[q_1,\; q_2,\; \dot{q}_1,\; \dot{q}_2]^{\top} \in \mathbb{R}^{4}
\]

where $q_1$ is the base-to-link1 joint angle and $q_2$ is the link1-to-link2
joint angle.

% ============================================================================
\section{Control Architecture}
% ============================================================================

\subsection{Cart-Pendulum Control Loop (OFC)}

The cart-pendulum control uses a biologically inspired OFC architecture with
four cascaded blocks:

\begin{center}
\textbf{LQR} $\xrightarrow{\bm{u}}$
\textbf{Muscle Dynamics} $\xrightarrow{\bm{F}}$
\textbf{ZFT Reference Mass} $\xrightarrow{\bm{s}_{\mathrm{ref}}}$
\textbf{Impedance Force} $\xrightarrow{\bm{F}_{\mathrm{imp}}}$
\textbf{Cart-Pendulum Plant}
\end{center}

\noindent\textbf{Important:} the muscle force $\bm{F}$ is fed into the
\emph{ZFT reference mass}, not directly into the impedance block.
The impedance block receives only $\bm{s}_{\mathrm{ref}}$ (position and
velocity of the reference mass) and $\bm{s}_{\mathrm{cart}}$ (actual cart
state). $\bm{F}$ reaches the impedance output only \emph{indirectly}, by
shaping how the reference mass moves.

\subsubsection*{Block 1 — Muscle Dynamics (Low-Pass Actuator Filter)}

The neural command $\bm{u} = [u_x, u_y]^{\top}$ is filtered through
first-order muscle dynamics:

\[
\boxed{
\dot{F}_i = \frac{-F_i + u_i}{\tau_m}, \quad i \in \{x, y\}
}
\]

\begin{itemize}
  \item \textbf{Input:} $\bm{u} \in \mathbb{R}^2$ (LQR neural command)
  \item \textbf{State:} $\bm{F} = [F_x, F_y]^{\top} \in \mathbb{R}^2$
  \item \textbf{Output:} $\bm{F} \in \mathbb{R}^2$ (muscle force)
  \item \textbf{Time constant:} $\tau_m = 0.03\;\mathrm{s}$ (human muscle)
\end{itemize}

This low-pass filter captures the finite bandwidth of biological muscles and
prevents instantaneous force jumps.

\subsubsection*{Block 2 — ZFT Reference Mass (Smooth Reference Trajectory)}

The muscle force $\bm{F}$ is the \emph{only} direct input from the muscle
block that propagates forward. It drives the ZFT reference mass ODE together
with a spring-damper coupling to the actual cart position. This is the
mechanism by which the neural command shapes the reference trajectory ---
$\bm{F}$ never enters the impedance block directly.

The ZFT dynamics are:

\[
\boxed{
M_{\mathrm{ref}}\,\ddot{\bm{p}}_{\mathrm{ref}}
= \bm{F} + K\,(\bm{p} - \bm{p}_{\mathrm{ref}})
         + D\,(\dot{\bm{p}} - \dot{\bm{p}}_{\mathrm{ref}})
}
\]

where $\bm{p} = [x, y]^{\top}$, $\bm{p}_{\mathrm{ref}} = [x_{\mathrm{ref}},\, y_{\mathrm{ref}}]^{\top}$.

\begin{itemize}
  \item \textbf{Inputs:}
        $\bm{s}_{\mathrm{cart}} = [x, y, \dot{x}, \dot{y}]^{\top} \in \mathbb{R}^4$
        and $\bm{F} \in \mathbb{R}^2$
  \item \textbf{State:}
        $\bm{s}_{\mathrm{ref}} = [x_{\mathrm{ref}}, y_{\mathrm{ref}},
        \dot{x}_{\mathrm{ref}}, \dot{y}_{\mathrm{ref}}]^{\top} \in \mathbb{R}^4$
  \item \textbf{Output:} $\bm{s}_{\mathrm{ref}} \in \mathbb{R}^4$
  \item \textbf{Parameters:} $M_{\mathrm{ref}} = 1.0\;\mathrm{kg}$,
        $K = 50\;\mathrm{N/m}$, $D = 10\;\mathrm{N\cdot s/m}$
\end{itemize}

The ZFT reference mass provides the ``efference copy'' of the intended
trajectory, decoupling perception from control.

\subsubsection*{Block 3 — Impedance Force (Spring-Damper Coupling)}

The impedance block computes the external force applied to the cart by
connecting it to the ZFT reference via a virtual spring-damper:

\[
\boxed{
\bm{F}_{\mathrm{imp}}
= K\,(\bm{p}_{\mathrm{ref}} - \bm{p})
+ D\,(\dot{\bm{p}}_{\mathrm{ref}} - \dot{\bm{p}})
}
\]

\begin{itemize}
  \item \textbf{Inputs:}
        $\bm{s}_{\mathrm{cart}} \in \mathbb{R}^4$ (actual cart pos.\ + vel.) and
        $\bm{s}_{\mathrm{ref}} \in \mathbb{R}^4$ (ZFT reference pos.\ + vel.).
        \textbf{The muscle force $\bm{F}$ is NOT a direct input here.}
  \item \textbf{Output:} $\bm{F}_{\mathrm{imp}} \in \mathbb{R}^2$ (cart actuation)
  \item \textbf{Parameters:} $K = 50\;\mathrm{N/m}$, $D = 10\;\mathrm{N\cdot s/m}$
\end{itemize}

$\bm{F}_{\mathrm{imp}}$ is applied directly to the cart's prismatic actuators.
The pendulum receives no direct torque input --- it is shaped purely by the
resulting cart motion.  The muscle force $\bm{F}$ influences $\bm{F}_{\mathrm{imp}}$
only \emph{indirectly}: by driving the ZFT ODE, it changes
$\bm{s}_{\mathrm{ref}}$, which in turn changes the spring-damper error seen
by the impedance block.

\subsubsection*{Block 4 — Finite-Horizon LQR Controller}

The LQR minimises the quadratic cost over a finite horizon $T$:

\[
\boxed{
J = \int_0^T \left[\bm{e}^{\top} Q\, \bm{e}
                 + \bm{u}^{\top} R\, \bm{u}\right] dt
  + \bm{e}(T)^{\top} Q_N\, \bm{e}(T)
}
\]

where $\bm{e} = \bm{x} - \bm{x}^{*}$ is the 14D state error with goal state

\[
\bm{x}^{*} =
[x^*, y^*, 0, 0, 0, 0, 0, 0, 0, 0, x^*, y^*, 0, 0]^{\top}
\]

The time-varying gain is obtained by solving the backward Riccati equation:

\[
-\dot{P}(t) = A^{\top}P + PA - PBR^{-1}B^{\top}P + Q,
\quad P(T) = Q_N
\]

yielding the control law:

\[
\bm{u}(t) = -R^{-1}B^{\top}P(t)\,\bm{e}(t)
\]

\paragraph{Cost matrices:}
\begin{align*}
Q &= \mathrm{diag}(100, 100, 500, 500, 10, 10, 100, 100, 0.1, 0.1, 1, 1, 0.1, 0.1) \\
Q_N &= 2\,Q \\
R &= \mathrm{diag}(1, 1)
\end{align*}

The system $(A, B)$ is obtained by linearising the 14D OFC dynamics around
the equilibrium (upright pendulum, zero velocity, zero force, cart at goal).

\subsection{Manipulator Control Loop}

The manipulator tracks the cart's 2D position using two cascaded blocks:

\begin{center}
\textbf{Cart Position} $\xrightarrow{\bm{p}_{\mathrm{cart}}}$
\textbf{IK Solver} $\xrightarrow{[q_d, \dot{q}_d]}$
\textbf{Computed-Torque Controller} $\xrightarrow{\bm{\tau}}$
\textbf{Manipulator Plant}
\end{center}

\subsubsection*{Block 5 — Inverse Kinematics (IK Solver)}

At every control step the desired cart position
$\bm{p}_{\mathrm{cart}} = [x, y]^{\top}$ is used as the EE target.
Drake's \texttt{InverseKinematics} solver finds joint angles
$\bm{q}_d = [q_{d,1}, q_{d,2}]^{\top}$ such that:

\[
\boxed{
f_{\mathrm{FK}}(\bm{q}_d) = \bm{p}_{\mathrm{cart}}, \quad
\|\bm{q}_d - \bm{q}_{\mathrm{prev}}\|_2 \to \min
}
\]

where $f_{\mathrm{FK}}$ is the forward-kinematics map of the manipulator and
warm-starting with the previous solution ensures continuity.
The desired joint velocities $\dot{\bm{q}}_d$ are computed numerically from
consecutive IK solutions.

\subsubsection*{Block 6 — Joint-Space Computed-Torque Controller}

A computed-torque law cancels the manipulator nonlinear dynamics and imposes
a linear PD law on the joint error:

\[
\boxed{
\bm{\tau} = M(\bm{q})\,\ddot{\bm{q}}_d
          + C(\bm{q}, \dot{\bm{q}})\,\dot{\bm{q}}
          + G(\bm{q})
          + K_p (\bm{q}_d - \bm{q})
          + K_d (\dot{\bm{q}}_d - \dot{\bm{q}})
}
\]

\begin{itemize}
  \item \textbf{Inputs:}
        desired state $[\bm{q}_d, \dot{\bm{q}}_d, \ddot{\bm{q}}_d]$
        and current state $[\bm{q}, \dot{\bm{q}}]$
  \item \textbf{Output:} $\bm{\tau} \in \mathbb{R}^2$ (joint torques)
  \item \textbf{Gains:} $K_p = 200\;\mathrm{N\cdot m/rad}$,
                        $K_d = 60\;\mathrm{N\cdot m\cdot s/rad}$,
                        $\tau_{\max} = 100\;\mathrm{N\cdot m}$
\end{itemize}

% ============================================================================
\section{Block Diagram}
% ============================================================================

\begin{figure}[H]
\centering
\includegraphics[width=\textwidth]{lqr_cart_manip_following.pdf}
\caption{Block diagram of the \texttt{lqr-applied-to-cart-manip-following-cart} mode.
         Solid arrows: signals in the primary control path.
         Dashed arrows: auxiliary feedback/feedforward signals.}
\label{fig:block_diagram}
\end{figure}

% ============================================================================
\section{Signal Flow Summary}
% ============================================================================

Table~\ref{tab:signals} summarises every signal connection in the diagram.

\begin{table}[H]
\centering
\caption{Signal connections in the control loop.}
\label{tab:signals}
\renewcommand{\arraystretch}{1.3}
\begin{tabular}{llll}
\toprule
\textbf{From} & \textbf{To} & \textbf{Signal} & \textbf{Dim.} \\
\midrule
LQR Controller       & Muscle Dynamics          & $\bm{u}$ (neural command)             & $\mathbb{R}^2$ \\
Muscle Dynamics      & ZFT Reference Mass       & $\bm{F}$ (muscle force)               & $\mathbb{R}^2$ \\
Muscle Dynamics      & State Mux (14D)          & $\bm{F}$ (for LQR state)             & $\mathbb{R}^2$ \\
Cart-Pendulum Plant  & ZFT Reference Mass       & $\bm{s}_{\mathrm{cart}}$              & $\mathbb{R}^4$ \\
Cart-Pendulum Plant  & Impedance Force          & $\bm{s}_{\mathrm{cart}}$              & $\mathbb{R}^4$ \\
ZFT Reference Mass   & Impedance Force          & $\bm{s}_{\mathrm{ref}}$ (\emph{only ZFT output goes to impedance, not} $\bm{F}$) & $\mathbb{R}^4$ \\
Impedance Force      & Cart-Pendulum Plant      & $\bm{F}_{\mathrm{imp}}$ (actuation) & $\mathbb{R}^2$ \\
Cart-Pendulum Plant  & State Mux                & $\bm{s}_{\mathrm{cart}}$ (8D)     & $\mathbb{R}^8$ \\
Muscle Dynamics      & State Mux                & $\bm{F}$                          & $\mathbb{R}^2$ \\
ZFT Reference Mass   & State Mux                & $\bm{s}_{\mathrm{ref}}$           & $\mathbb{R}^4$ \\
State Mux            & ZOH (state hold)         & $\bm{x}$ (14D)                    & $\mathbb{R}^{14}$ \\
ZOH                  & LQR Controller           & $\bm{x}$ (14D, held)              & $\mathbb{R}^{14}$ \\
\midrule
Cart-Pendulum Plant  & IK Solver                & $[x, y, \dot{x}, \dot{y}]$       & $\mathbb{R}^4$ \\
IK Solver            & Computed-Torque Ctrl.    & $[\bm{q}_d, \dot{\bm{q}}_d, \ddot{\bm{q}}_d]$ & $\mathbb{R}^6$ \\
Manipulator Plant    & Computed-Torque Ctrl.    & $[\bm{q}, \dot{\bm{q}}]$         & $\mathbb{R}^4$ \\
Computed-Torque Ctrl.& Manipulator Plant        & $\bm{\tau}$ (joint torques)       & $\mathbb{R}^2$ \\
\bottomrule
\end{tabular}
\end{table}

% ============================================================================
\section{Linearised System}
% ============================================================================

The LQR is designed on the linearised 14D OFC system.
Let $\bm{x}$ be the full state and $\bm{u}$ the neural command.
Around the equilibrium $(\bm{x}^*, \bm{u}^* = \bm{0})$ the dynamics are:

\[
\dot{\bm{x}} = A\,\bm{x} + B\,\bm{u}
\]

The $A \in \mathbb{R}^{14 \times 14}$ and $B \in \mathbb{R}^{14 \times 2}$
matrices encode:
\begin{itemize}
  \item \textbf{Cart-Pendulum linearisation:} obtained from the Drake
        \texttt{Linearize()} helper applied to the 8D \texttt{MultibodyPlant}
        at the upright equilibrium.
  \item \textbf{Muscle augmentation:} the muscle dynamics
        $\dot{F}_i = (-F_i + u_i)/\tau_m$ extend the $A$ and $B$ matrices by
        two rows/columns each.
  \item \textbf{ZFT augmentation:} the reference-mass dynamics extend the
        state by four more rows.
\end{itemize}

The discrete-time Riccati equation is solved backward in time from $t = T$ to
$t = 0$ at timestep $\Delta t = 10\;\mathrm{ms}$ to obtain $P(t)$, then the
gain schedule $K(t) = R^{-1}B^{\top}P(t)$ is stored and interpolated online.

% ============================================================================
\section{Simulation Parameters}
% ============================================================================

\begin{table}[H]
\centering
\caption{Default simulation parameters.}
\renewcommand{\arraystretch}{1.3}
\begin{tabular}{lll}
\toprule
\textbf{Parameter} & \textbf{Value} & \textbf{Description} \\
\midrule
Cart mass           & $3.0\;\mathrm{kg}$  & \\
Pendulum mass       & $0.3\;\mathrm{kg}$  & \\
Pendulum length     & $0.25\;\mathrm{m}$  & \\
Cart damping        & $0.5\;\mathrm{N\cdot s/m}$ & \\
Pendulum damping    & $0.1\;\mathrm{N\cdot s/rad}$ & \\
$\tau_m$            & $0.03\;\mathrm{s}$   & Muscle time constant \\
$M_{\mathrm{ref}}$  & $1.0\;\mathrm{kg}$   & ZFT reference mass \\
$K$                 & $50\;\mathrm{N/m}$   & Impedance stiffness \\
$D$                 & $10\;\mathrm{N\cdot s/m}$ & Impedance damping \\
$K_p$               & $200\;\mathrm{N\cdot m/rad}$ & Joint PD: position gain \\
$K_d$               & $60\;\mathrm{N\cdot m\cdot s/rad}$ & Joint PD: velocity gain \\
LQR horizon $T$     & $10.0\;\mathrm{s}$   & \\
$\Delta t$ (LQR)    & $10\;\mathrm{ms}$    & Riccati integration step \\
Simulation $\Delta t$ & $1\;\mathrm{ms}$   & Drake integrator step \\
\bottomrule
\end{tabular}
\end{table}

% ============================================================================
\section{Implementation Notes}
% ============================================================================

\begin{itemize}
  \item A \textbf{Zero-Order Hold (ZOH)} at $10\;\mathrm{ms}$ is inserted
        between the 14D state multiplexer and the LQR input to break the
        algebraic loop introduced by directly feeding the LQR output (through
        muscle and impedance) back to itself via the ZFT state.

  \item The \textbf{IK solver} is warm-started with the previous joint
        solution at every step to ensure smooth, continuous joint trajectories
        and avoid large discontinuous jumps.

  \item All systems are \textbf{Drake LeafSystems} wired inside a single
        \texttt{DiagramBuilder}, which guarantees correct simulation semantics
        (context-based state isolation, publication order, etc.).

  \item The cart-pendulum and manipulator share the \textbf{same
        \texttt{MultibodyPlant}} but have separate \texttt{ModelInstance}
        handles, enabling independent state access and actuation while
        capturing any physical coupling.
\end{itemize}

% ============================================================================
\newpage
\section{Equation Flow Summary}
% ============================================================================

This section presents the complete signal chain as a single, self-contained
narrative — from sensor reading to actuator output — showing exactly how each
equation feeds into the next at every simulation step.

\subsection*{Step-by-step signal chain}

\noindent At every simulation step $t$, the following sequence executes:

\bigskip
\noindent\textbf{(1) Read plant state.}
The cart-pendulum plant outputs its 8D state:
\[
\bm{s}_{\mathrm{cart}}(t)
= \bigl[\,x,\;y,\;\alpha,\;\beta,\;\dot{x},\;\dot{y},\;\dot{\alpha},\;\dot{\beta}\,\bigr]^{\!\top}
\in\mathbb{R}^8
\]
A demultiplexer splits this into four 2-vectors
$[\bm{p}, \bm{\theta}, \dot{\bm{p}}, \dot{\bm{\theta}}]$,
and a multiplexer re-assembles the \emph{cart-only} sub-state:
\[
\bm{s}_{\mathrm{c}}(t) = [\,x,\;y,\;\dot{x},\;\dot{y}\,]^{\!\top} \in\mathbb{R}^4
\]

\bigskip
\noindent\textbf{(2) Muscle dynamics — low-pass filter on neural command.}
The LQR command $\bm{u} \in\mathbb{R}^2$ (from the previous ZOH slot) drives
first-order actuator dynamics:
\[
\dot{\bm{F}}(t) = \frac{-\bm{F}(t) + \bm{u}(t)}{\tau_m}
\quad\Longrightarrow\quad
\bm{F}(t) \in\mathbb{R}^2
\]
The muscle state $\bm{F}$ is (a) fed to the \textbf{ZFT reference mass only}
(not to the impedance block directly), and (b) stacked into the 14D state
for the LQR. $\bm{F}$ reaches the impedance output only indirectly, by
shaping how the reference mass moves.

\bigskip
\noindent\textbf{(3) ZFT reference mass — smooth reference trajectory.}
The virtual reference mass is driven by muscle force and the error to the cart:
\[
\ddot{\bm{p}}_{\mathrm{ref}}(t)
= \frac{1}{M_{\mathrm{ref}}}\!\left[
    \bm{F}(t)
    + K\bigl(\bm{p}(t) - \bm{p}_{\mathrm{ref}}(t)\bigr)
    + D\bigl(\dot{\bm{p}}(t) - \dot{\bm{p}}_{\mathrm{ref}}(t)\bigr)
  \right]
\]
Integrating gives:
\[
\bm{s}_{\mathrm{ref}}(t) = [\,x_{\mathrm{ref}},\;y_{\mathrm{ref}},\;\dot{x}_{\mathrm{ref}},\;\dot{y}_{\mathrm{ref}}\,]^{\!\top}
\in\mathbb{R}^4
\]

\bigskip
\noindent\textbf{(4) Impedance force --- virtual spring-damper actuation.}
Using only $\bm{s}_{\mathrm{c}}$ (actual cart state) and
$\bm{s}_{\mathrm{ref}}$ (ZFT output), the impedance block
computes the external force applied to the cart:
\[
\bm{F}_{\mathrm{imp}}(t)
= K\bigl(\bm{p}_{\mathrm{ref}} - \bm{p}\bigr)
+ D\bigl(\dot{\bm{p}}_{\mathrm{ref}} - \dot{\bm{p}}\bigr)
\in\mathbb{R}^2
\]
\textbf{Note:} $\bm{F}$ (muscle force) is \emph{not} a direct input here.
It enters this equation only indirectly through $\bm{p}_{\mathrm{ref}}$,
which was shaped by the ZFT ODE in Step~3.
This is the \emph{sole actuation input} to the cart-pendulum plant.
The pendulum angles $\alpha,\beta$ are shaped indirectly by the resulting
cart acceleration.

\bigskip
\noindent\textbf{(5) Full 14D state assembly.}
All sub-states are concatenated:
\[
\bm{x}(t) =
\underbrace{[\,x,\;y,\;\alpha,\;\beta\,]}_{\text{cart pos.}}
\oplus
\underbrace{[\,\dot{x},\;\dot{y},\;\dot{\alpha},\;\dot{\beta}\,]}_{\text{cart vel.}}
\oplus
\underbrace{[\,F_x,\;F_y\,]}_{\text{muscle}}
\oplus
\underbrace{[\,x_{\mathrm{ref}},\;y_{\mathrm{ref}},\;\dot{x}_{\mathrm{ref}},\;\dot{y}_{\mathrm{ref}}\,]}_{\text{ZFT ref.}}
\in\mathbb{R}^{14}
\]
A Zero-Order Hold (ZOH, $\Delta = 10\;\mathrm{ms}$) samples this to
$\hat{\bm{x}}(t)$ before passing it to the LQR.

\bigskip
\noindent\textbf{(6) Finite-horizon LQR — optimal neural command.}
The state error $\bm{e}(t) = \hat{\bm{x}}(t) - \bm{x}^{*}$ enters the
time-varying gain computed by the backward Riccati pass:
\[
\bm{u}(t) = -\underbrace{R^{-1}B^{\!\top}P(t)}_{K(t)}\,\bm{e}(t)
\in\mathbb{R}^2
\]
This $\bm{u}(t)$ closes the outer loop back to \textbf{Step 2} (muscle input)
at the next time step.

\bigskip
\noindent\textbf{(7) IK solver — cart position to joint angles.}
In parallel, the cart position $[x,y]$ is used as the EE target for the
manipulator's inverse kinematics:
\[
\bm{q}_d(t) = f_{\mathrm{IK}}\!\bigl([x(t),\,y(t)],\;\bm{q}_d(t{-}\Delta t)\bigr)
\in\mathbb{R}^2
\]
Desired joint velocity is approximated numerically:
$\dot{\bm{q}}_d = (\bm{q}_d(t) - \bm{q}_d(t{-}\Delta t))/\Delta t$.

\bigskip
\noindent\textbf{(8) Computed-torque controller — joint torques.}
The nonlinear manipulator dynamics are cancelled and a PD law is imposed:
\[
\bm{\tau}(t)
= M(\bm{q})\ddot{\bm{q}}_d
+ C(\bm{q},\dot{\bm{q}})\dot{\bm{q}}
+ G(\bm{q})
+ K_p(\bm{q}_d - \bm{q})
+ K_d(\dot{\bm{q}}_d - \dot{\bm{q}})
\in\mathbb{R}^2
\]
$\bm{\tau}$ is applied directly to the manipulator joints,
driving the end-effector to follow the cart trajectory.

\subsection*{Compact causal chain}

\noindent The complete equation flow can be read as a single chain of
dependencies:

\[
\underbrace{\bm{s}_{\mathrm{cart}}}_{\text{plant}}
\;\xrightarrow{\text{demux/mux}}\;
\bm{s}_{\mathrm{c}}
\;\xrightarrow[\text{ZFT}]{\bm{F}}\;
\bm{s}_{\mathrm{ref}}
\;\xrightarrow{\text{imp.}}\;
\bm{F}_{\mathrm{imp}}
\;\longrightarrow\;
\text{cart plant}
\]
\[
\bm{s}_{\mathrm{cart}} \oplus \bm{F} \oplus \bm{s}_{\mathrm{ref}}
\;\xrightarrow{\text{mux}}\;
\bm{x}_{14}
\;\xrightarrow{\text{ZOH}}\;
\hat{\bm{x}}
\;\xrightarrow{\text{LQR}}\;
\bm{u}
\;\xrightarrow{\tau_m}\;
\bm{F}
\quad\text{(closes loop)}
\]
\[
\bm{s}_{\mathrm{c}}
\;\xrightarrow{\text{IK}}\;
\bm{q}_d, \dot{\bm{q}}_d
\;\xrightarrow{\text{CTC}}\;
\bm{\tau}
\;\longrightarrow\;
\text{manip.\ plant}
\]

\noindent Each arrow represents a Drake system connection.
The cart-pendulum loop (top two lines) and the manipulator loop (bottom line)
run concurrently inside the same Drake \texttt{Diagram} at every integration
step.

\end{document}
